For a given DFT frequency series $\tilde x[l]$, the WDM transform represents the data through the WDM coefficients $w_{nm}$
\begin{align}\label{wnm_def}
w_{nm}&=\sum_{l=-N/2}^{N/2-1}\tilde x[l]\tilde g^*_{nm}[l]\,,\codelink{https://github.com/pywavelet/wdm_transform/blob/v0.5.0/src/wdm_transform/transforms/xp_backend.py\#L225}
\end{align}
where $n$ corresponds to the time index and $m$ corresponds to the
frequency-channel.

The Wilson-basis elements $\tilde{g}$ with Daubechies-Meyer frequency-domain window $\tilde{\varphi}$ are defined in closed form in the discrete frequency domain below. For DC ($m=0$), interior channels ($0 < m < N_f$), and
Nyquist ($m=N_f$) we have that:
\begin{widetext}
\begin{align}
\label{eq:gnm_full_discrete}
\text{\textbf{Wilson Basis:}}\quad \tilde{g}_{nm}[l] &=
\dfrac{1}{\sqrt{2}}\begin{cases}
e^{-4\pi i n l/N_t}\, \tilde{\varphi}[l], & m=0, \\[6pt]
e^{-2\pi i n l/N_t}\Big(C_{nm}\,\tilde{\varphi}[l - mN_t/2] + C_{nm}^*\,\tilde{\varphi}[l + mN_t/2]\Big), & 0 < m < N_f, \\[6pt]
e^{-4\pi i n l/N_t}\Big(\tilde{\varphi}[l - N/2] + \tilde{\varphi}[l + N/2]\Big), & m = N_f, 
\end{cases}
\codelink{https://github.com/pywavelet/wdm_transform/blob/v0.5.0/src/wdm_transform/windows.py\#L129}\\
\text{\textbf{Daubechies--Meyer:}} \quad  \tilde\varphi[l]&=\sqrt{\frac{2}{N_t}}\begin{cases}
1 & |l_r|<A \\
\cos \left[\frac{\pi}{2}\,\left(
\frac{|l_r|-A}{B}\right)\right] & A \leq|l_r| < A+B\,, \\
0 & \text{otherwise}\,,
\end{cases} \qquad l_r =l\frac{2}{N_t}\,,\codelink{https://github.com/pywavelet/wdm_transform/blob/v0.5.0/src/wdm_transform/windows.py\#L71} \label{eq:DM_window_discrete}\\
  C_{nm}&=e^{\frac{i\pi}{4}(1-(-1)^{n+m})} =
  \begin{cases} 1 & n+m \text{ even,} \\ i & n+m \text{ odd.}\label{eq:c_nm_def}
\end{cases}\codelink{https://github.com/pywavelet/wdm_transform/blob/v0.5.0/src/wdm_transform/windows.py\#L58}
\end{align}
\begin{equation}\label{eq:A_B_condition}
B = 1 - 2A \,, \quad \text{for}\ A\in(0,1/2)\,. 
\end{equation}
\end{widetext}
Equation~\eqref{eq:A_B_condition} gives freedom to the user to define the properties of the Daubechies--Meyer window defined in Eq.~\eqref{eq:DM_window_discrete}.

From this point onwards, we will provide an interpretation for each of the equations \eqref{eq:gnm_full_discrete} -- \eqref{eq:A_B_condition}. We begin by discussing the DC, Interior, Nyquist split of the basis functions $\tilde{g}$, an important feature of our codebase, and the practical and theoretical consequences of the Daubechies--Meyer window $\tilde{\varphi}$.

% This orthogonality is the key reason that the WDM formalism is adept for treating stationary noise processes for reasons that will be discussed later.

From Eq.\eqref{eq:c_nm_def}, the interior channels carry an alternating phase factor $C_{nm}$, while the DC and Nyquist edge channels have a doubled time-shift exponent and no $C_{nm}$ modulation. The phase factor $C_{nm}$ is the
``\emph{Wilson}'' ingredient of the WDM construction: it combines the
positive- and negative-frequency windowed Fourier atoms centered at $mN_t/2$ into a single real, orthogonal pair, which is what
makes the interior basis real-valued and Wilson-orthogonal, to be discussed later in this section, in the
discrete setting. 

These edge channels $m = 0$ and $m = N_f$ in \eqref{eq:gnm_full_discrete} each contribute $N_t$ real degrees of freedom to
the packed grid: the interior block stores $N_t(N_f{-}1)$ real
coefficients via the $C_{nm}$-modulated pairing of the
$\pm mN_t/2$ atoms (the Wilson-basis construction), while the DC
and Nyquist channels each contribute $N_t$ unpaired real coefficients.
The total $N_t(N_f{+}1)$ stored entries therefore represent the same
$N=N_tN_f$ time-domain degrees of freedom, with the edge channels
being exactly $N_t$ real numbers~\cite{Necula_2012,Cornish:2020odn}. 
We refer the reader to Fig.~\ref{fig:wdm_transform} for a schematic of the WDM transform and the packed coefficient layout.  % 

The window $\tilde\varphi$ is the ``\emph{Daubechies--Meyer}''
ingredient of the WDM construction: a smooth, compactly supported
frequency window of the Meyer family, whose taper region overlaps
between adjacent channels in a way that yields exact partition-of-unity
reconstruction. The Meyer family is parametrized by an integer order
$d \geq 1$ that controls the smoothness of the taper: the cosine
argument $(|l_r|-A)/B$ in Eq.~\eqref{eq:DM_window_discrete} is replaced
by the normalized incomplete beta function
\begin{equation}\label{eq:nu_d_def}
  \nu_d(y) = \frac{\int_0^y t^{d-1}(1-t)^{d-1}\,\mathrm{d}t}
                  {\int_0^1 t^{d-1}(1-t)^{d-1}\,\mathrm{d}t}\,,
  \qquad y = \frac{|l_r|-A}{B}\,,
\end{equation}
which satisfies $\nu_d(y) + \nu_d(1-y) = 1$ and steepens the roll-off
as $d$ increases. Setting $d=1$ gives $\nu_1(y)=y$ and recovers the
cosine taper of Eq.~\eqref{eq:DM_window_discrete}. In the present
implementation, and in this work, we use this
$d=1$ cosine-tapered member of the family. We do stress that the parameter $d$ is retained in the Application Programming Interface (API) only for
possible future
generalizations.\footnote{Cornish~\cite{Cornish:2020odn} uses $d=4$,
  which yields a flatter passband and steeper roll-off at the expense
  of a wider time-domain support (the window function must be truncated
  at $\pm q\Delta T$ with $q=16$). The $d=1$ cosine taper used here has
  a gentler roll-off but is more compact in time, requiring a smaller
  $q$. In practice, the choice of $d$ trades spectral leakage
  suppression against temporal compactness of the window; for the
  applications considered in this paper the $d=1$ window provides
  adequate frequency separation while keeping the implementation simple
and the time-domain filter short.} By construction, $\tilde\varphi$ has a limited frequency bandwidth
which is controlled by the parameter $A\in(0,1/2)$ in Eq.\eqref{eq:A_B_condition}. 
The frequency localization of the WDM atoms is illustrated in Fig.~\ref{fig:wdm_transform}(b), where the same two highlighted cells are shown in frequency space.
At each wavelet frequency bin $l$, one sees that the window overlaps with adjacent frequency bins.  

For our implementation we use the default $A =1/3$ and
hence $B = 1/3$, so the window is non-zero for normalised frequencies
$\leq 2/3$. This is a default, not a hard-coded
constraint: $A$ is exposed in \texttt{wdm\_transform} as the
parameter \texttt{a} of the \texttt{WDM} class (with $B=1-2A$ derived
automatically), and any choice with $A\in(0,1/2)$ is admissible.%

A final and essential property of the WDM formalism is the principle of orthogonality. The Wilson basis $\tilde{g}$ defined by Eq.~\eqref{eq:gnm_full_discrete} combined with the Daubechies--Meyer window Eq.~\eqref{eq:DM_window_discrete} with associated constants~\eqref{eq:c_nm_def} provide an orthogonal basis\footnote{Our definitions of the Wilson basis functions \eqref{eq:gnm_full_discrete} carefully include the $m=0$ and $m = N_f$ in order for orthogonality to hold.}, with both of the following equalities proved in Appendix~\ref{app:normality_condition} using Eq.~\eqref{eq:DM_window_discrete} and Eq.~\eqref{eq:c_nm_def}:%
\begin{widetext}
\begin{align}\label{eq:orthonormality}
\sum_{l=-N/2}^{N/2-1}\tilde{g}^*_{nm}[l]\tilde{g}_{pq}[l]&=\delta_{m,q}\begin{cases}
    \frac{1}{2}\left(\delta_{n,p}+\delta_{n,p\pm \frac{N_t}{2}}\right) \qquad m\in\{0,N_f\}\\
    \delta_{n,p} \qquad 0<m<N_f
\end{cases}\nonumber\\
\sum_{n=0}^{N_t-1}\sum_{m=0}^{N_f}\tilde{g}^*_{nm}[l]\tilde{g}_{nm}[l']&=\delta_{ll'}\;,
\end{align}
\end{widetext}

The basis functions here are represented in the frequency domain and it is easily proved they are orthogonal in the time domain.
Orthogonality is one of the key reasons that the WDM formalism is adept for gravitational-wave data analysis, which will be discussed at a later stage of this work.



These are the main ingredients behind the packed WDM grid. A single window shape is shifted across
frequency to define the channel index $m$, and each shifted window is then modulated in time to define the atom indexed by $n$. This is why one WDM coefficient can be read as the strength of one localized time-frequency pattern, rather than as the amplitude of a globally supported sinusoid. 

More concretely, a single coefficient $w_{nm}$ answers the question:
``how much does the data resemble the atom centered at time bin $n$
and frequency bin $m$?'' A large magnitude means that the data
contains a feature with roughly the same time localization and central
frequency. 
Reading the WDM plane coefficient-by-coefficient in this way is often much more
intuitive than interpreting a global Fourier series, because the
index pair $(n,m)$ already carries an approximate time-frequency location.

In the next subsections, we will provide formulae for the $w_{nm}$ coefficients for both the interior, DC and Nyquist channels separately. In our final subsection, we investigate the orthogonal nature of the Wilson-basis.  

\subsubsection{Interior channels \texorpdfstring{$(0 < m < N_f)$}{(0 < m < N_f)}}

Starting from Eq.~\eqref{wnm_def} with the interior-channel basis
element from Eq.~\eqref{eq:gnm_full_discrete}, the forward transform can be
evaluated in the Fourier domain using
\begin{align}
  w_{nm} = \sum_{l=-N/2}^{N/2-1} \tilde{x}[l]\,\tilde{g}_{nm}^*[l]\,,
\end{align}
For a real-valued input signal, $\tilde{x}[-l] = \tilde{x}^*[l]$, and the real-valued window satisfies $\tilde{\varphi}[-l] = \tilde{\varphi}[l]$.
Splitting the sum into positive and negative frequency halves and
using these symmetries, the negative-frequency contribution is the
complex conjugate of the positive-frequency contribution (see
Appendix~\ref{app:forward_derivation} for details), giving
%
\begin{equation}\label{eq:wnm_interior}
\begin{split}
  w_{nm} = \sqrt{2}&\,(-1)^{nm}\,\Re\Bigl[C_{nm}^*\\
  &\times\sum_{l'=-N_t/2}^{N_t/2-1}\tilde{x}[l'{+}mN_t/2]\,e^{2\pi i l'n/N_t}\,\tilde{\varphi}[l']\Bigr],\codelink{https://github.com/pywavelet/wdm_transform/blob/v0.5.0/src/wdm_transform/transforms/xp_backend.py\#L92-L102}
\end{split}
\end{equation}
where $l' = l - mN_t/2$ and $\Delta f = 1/T$.
The expression above can be written compactly as
%
\begin{align}\label{eq:wnm_compact}
  w_{nm} = \sqrt{2}\,(-1)^{nm}\,\Re\,C_{nm}^*\,\tilde x_m[n]\,,
\end{align}
%
where
%
\begin{align}\label{eq:xm_def}
  \tilde x_m[n] = \sum_{l=-N_t/2}^{N_t/2-1} e^{2\pi i l n/N_t}\,\tilde x[l +
  mN_t/2]\,\tilde{\varphi}[l],
\end{align}
%
The quantity
$\tilde x_m[n]$ can be computed using a length-$N_t$ FFT for each channel
$m$, so the total cost of the interior channels is $N_f \times
\mathcal{O}(N_t\log N_t)$.

\subsubsection{DC channel \texorpdfstring{$(m = 0)$}{(m = 0)}}

The DC edge channel has a doubled time-shift exponent and no $C_{nm}$
factor in the basis definition~\eqref{eq:gnm_full_discrete}. Following the
same Fourier-domain evaluation yields
\begin{align}\label{eq:wn0}
  w_{n0} = \sqrt{2\,}\Re\left[\sum_{l=1}^{N_t/2-1}\tilde{x}[l]\,e^{4\pi i l
  n/N_t}\,\tilde{\varphi}[l]\right] + \frac{1}{\sqrt{2}}\tilde{x}[0]\,\tilde{\varphi}[0].\codelink{https://github.com/pywavelet/wdm_transform/blob/v0.5.0/src/wdm_transform/transforms/xp_backend.py\#L80-L90}
\end{align}
Note the factor of $4\pi$ in the exponent (versus $2\pi$ for interior
channels): the DC basis element oscillates at twice the rate in the
time index because it sees both positive and negative frequency
copies of the window centered at $f=0$.

\subsubsection{Nyquist channel \texorpdfstring{$(m = N_f)$}{(m = N_f)}}

The Nyquist edge channel is handled analogously to the DC channel but
centered at the Nyquist frequency $f_{\rm max}$. Its structure
mirrors the DC case with appropriate frequency shifts. In practice,
the Nyquist channel is computed using the same windowed-FFT machinery
as the interior channels, with the frequency shift set to $mN_t/2 = N/2$.
%
\begin{align}\label{eq:wnNf}
  w_{nN_f} &= \sqrt{2}\,\Re\!\left[\sum_{l=N/2-N_t/2}^{N/2-1}
    \tilde{x}[l]\,e^{4\pi i l n/N_t}\,\tilde{\varphi}[l-N/2]\right]
    \nonumber\\
  &\quad + \frac{1}{\sqrt{2}}\,\tilde{x}[N/2]\,\tilde{\varphi}[0]\,.\codelink{https://github.com/pywavelet/wdm_transform/blob/v0.5.0/src/wdm_transform/transforms/xp_backend.py\#L104-L114}
\end{align}

Collecting the results of each of the $w_{nm}$ coefficients above, the forward transform reads
\begin{widetext}
\begin{align}
\label{eq:forward_transform_discrete}
w_{nm} = \sqrt{2}(-1)^{mn}\Re\,
\begin{cases}
\,\displaystyle\sum_{l=1}^{N_t/2 - 1} \tilde{x}[l]\, \tilde{\varphi}[l]\, e^{4\pi i n l/N_t}\,
+ \dfrac{1}{2}\, \tilde{x}[0]\, \tilde{\varphi}[0],
& m = 0, \\[12pt]
 C_{nm}^{*} \displaystyle\sum_{l' = -N_t/2}^{N_t/2 - 1} \tilde{x}[l' + mN_t/2]\, \tilde{\varphi}[l']\, e^{2\pi i l' n/N_t}\,,
& 1 \le m \le N_f - 1, \\[12pt]
\displaystyle\sum_{l = N/2 - N_t/2}^{N/2 - 1} \tilde{x}[l]\, \tilde{\varphi}[l - N/2]\, e^{4\pi i n l/N_t}
+ \dfrac{1}{2}\, \tilde{x}[N/2]\, \tilde{\varphi}[0],
& m = N_f,
\end{cases}
\codelink{https://github.com/pywavelet/wdm_transform/blob/v0.5.0/src/wdm_transform/transforms/xp_backend.py\#L65-L116}
\end{align}
\end{widetext}
holding true for $N_f \in 2\mathbb{N}$. 


